\section{Exact representation frontiers}\label{sec:theory-v1}

We give set-theoretic results first and state their limits explicitly.  The maximal sound set below
is the complement-of-encoded-unsafe-set construction already used by
\citet{lutkus2025latent}; safe/unsafe collision and qualitative conservatism are therefore not our
novelty claim.  Our purpose is to fix the endpoint, separate observation from encoder ambiguity,
and extend the analysis to common safe actions.  All suprema and infima are extended-real, with
$\sup(\varnothing\cup\{0\})=0$ and $\inf\varnothing=+\infty$.

\subsection{Static certificates and endpoint conventions}

Let
\[
 \mathcal C=\{x\in\K:h(x)\geq0\},\qquad
 \mathcal D=\{x\in\K:h(x)<0\},
\]
and let $R:\K\to\Z$ be deterministic.  A latent set $S\subseteq\Z$ is
\emph{sound} if $R^{-1}(S)\subseteq\mathcal C$.  For $m\geq0$, it is closed
$m$-complete if $\{x:h(x)\geq m\}\subseteq R^{-1}(S)$ and open $m$-complete if
$\{x:h(x)>m\}\subseteq R^{-1}(S)$.  Define
\begin{align*}
 \Gamma_R
 &=\{h(x):h(x)\geq0,\ \exists x'\in\mathcal D\text{ with }R(x')=R(x)\},\\
 \epssep(R)&=\sup(\Gamma_R\cup\{0\}).
\end{align*}

\begin{theorem}[Exact existence equivalence]\label{thm:exact-existence}
Let $S_{\max}=\Z\setminus R(\mathcal D)$.  Then $S$ is sound if and only if
$S\subseteq S_{\max}$.  Moreover, a sound closed $m$-complete set exists if and only if
\begin{equation}
 \Gamma_R\cap[m,+\infty)=\varnothing,\label{eq:closed-completeness}
\end{equation}
whereas a sound open $m$-complete set exists if and only if
\begin{equation}
 \Gamma_R\cap(m,+\infty)=\varnothing.\label{eq:open-completeness}
\end{equation}
Whenever the relevant condition holds, $S_{\max}$ attains it.
\end{theorem}

\begin{proof}
If a sound $S$ contained $z\in R(\mathcal D)$, an unsafe preimage of $z$ would belong to
$R^{-1}(S)$, a contradiction.  Conversely, a preimage of a point outside $R(\mathcal D)$ cannot
be unsafe, so every subset of $S_{\max}$ is sound.  Since every sound set is contained in
$S_{\max}$, closed $m$-completeness is possible precisely when no $x$ with $h(x)\geq m$ has a
code in $R(\mathcal D)$.  This is exactly \eqref{eq:closed-completeness}.  Replacing $\geq m$ by
$>m$ proves \eqref{eq:open-completeness}.
\end{proof}

\begin{proposition}[Infimum completeness threshold]\label{prop:threshold}
Define
\begin{align*}
 \tau_{\geq}(R)
 &=\inf\{m\geq0:\exists\text{ sound closed $m$-complete }S\},\\
 \tau_{>}(R)
 &=\inf\{m\geq0:\exists\text{ sound open $m$-complete }S\}.
\end{align*}
Then
\[
 \tau_{\geq}(R)=\tau_{>}(R)=\epssep(R).
\]
For finite $\epssep$, open completeness is attainable at $m=\epssep$; closed completeness is
attainable there if and only if $\epssep\notin\Gamma_R$.  If $\epssep=+\infty$, no finite
margin is attainable.
\end{proposition}

\begin{proof}
For every finite $m>\epssep$, \eqref{eq:closed-completeness} holds.  For
$0\leq m<\epssep$, the definition of supremum gives an element of $\Gamma_R$ strictly larger than
$m$, so neither completeness condition holds.  At a finite endpoint no element exceeds
$\epssep$, proving open attainability, while \eqref{eq:closed-completeness} holds exactly when the
supremum is not itself in $\Gamma_R$.  If the supremum is infinite, $\Gamma_R$ is unbounded.
Taking infima proves the claim.
\end{proof}

\begin{remark}[Boundary and endpoint]
If a safe boundary state with $h=0$ shares a code with an unsafe state, then $\epssep=0$, yet
closed zero-completeness is impossible.  Thus Proposition~\ref{prop:threshold} is an infimum
statement, not an endpoint-achievement statement.  Likewise, when conflicting margins approach
one from below but no margin-one state is conflicted, closed one-completeness is attainable even
though $\epssep=1$.
\end{remark}

\subsection{Deterministic data processing}

\begin{theorem}[Observation floor]\label{thm:data-processing}
If $R_2=T\circ R_1$ for deterministic maps $R_1$ and $T$, then
\[
 \Gamma_{R_1}\subseteq\Gamma_{R_2},\qquad
 \epssep(R_1)\leq\epssep(R_2),\qquad
 \tau_{\geq}(R_1)\leq\tau_{\geq}(R_2).
\]
In particular, for deterministic observation $O$ and encoder $E$,
\[
 \epssep(O)\leq\epssep(E\circ O).
\]
\end{theorem}

\begin{proof}
Equality under $R_1$ implies equality after applying $T$, so every conflicting pair for $R_1$ is
also conflicting for $R_2$.  Set inclusion, the supremum inequality, and
Proposition~\ref{prop:threshold} give the three conclusions.
\end{proof}

This theorem concerns exact fibers.  A fixed raw latent radius has no such invariant law: an
injective rescaling preserves exact fibers while moving pairs into or out of the radius.  A history
encoder can beat a memoryless observation floor only by receiving a richer input history or
belief, rather than by post-processing the same observation.

\subsection{Safe-action fibers}

Let $F_z^R=\{x:R(x)=z\}$ and let $\Usafe(x)$ denote the actions satisfying a declared audited
condition.  Define the common-action set
\[
 \mathcal A_R(z)=\bigcap_{x\in F_z^R}\Usafe(x).
\]
A fiber is action-conflicted if every $\Usafe(x)$ is nonempty but $\mathcal A_R(z)=\varnothing$.

\begin{theorem}[Deterministic action conflict]\label{thm:det-action}
A deterministic memoryless decision $\pi(z)$ is safe for every $x\in F_z^R$ if and only if
$\pi(z)\in\mathcal A_R(z)$.  Consequently, such a uniformly safe decision exists at $z$ if and
only if $\mathcal A_R(z)\neq\varnothing$.  Hence an action-conflicted fiber rules out every
uniformly safe deterministic memoryless latent policy.
\end{theorem}

\begin{proof}
The decision is safe for every compatible state exactly when the single action $\pi(z)$ belongs
to their safe-action intersection.
\end{proof}

Let $(\mathcal U,\Sigma_{\mathcal U})$ be a measurable action space.  For a randomized decision $\mu_z$,
assume that every $\Usafe(x)$ and $\mathcal A_R(z)$ belongs to $\Sigma_{\mathcal U}$.  Distinguish strong safety
$\mu_z(\mathcal A_R(z))=1$ from pointwise almost-sure safety
$\mu_z(\Usafe(x))=1$ for every fixed $x\in F_z^R$.

\begin{theorem}[Randomized action conflict]\label{thm:rand-action}
A strongly safe randomized decision exists if and only if $\mathcal A_R(z)$ is nonempty.  The same
equivalence holds under pointwise almost-sure safety if either (i) the fiber is countable, or
(ii) the action space is compact metric, $\mu_z$ is a Borel probability measure, and every
$\Usafe(x)$ is closed.
\end{theorem}

\begin{proof}
A probability measure cannot give mass one to the empty common set, while a common action induces
a safe Dirac measure.  For a countable fiber, the intersection of the probability-one safe-action
events still has probability one.  In case (ii), every finite intersection has probability one
and is nonempty.  The closed safe-action sets therefore have the finite-intersection property, and
compactness makes their full intersection nonempty.  The converse in both cases is again a Dirac
measure.
\end{proof}

Some qualification in Theorem~\ref{thm:rand-action} is necessary.  With action space $[0,1]$, an
uncountable fiber indexed by $x\in[0,1]$, and $\Usafe(x)=[0,1]\setminus\{x\}$, the common set is
empty but the uniform distribution is pointwise almost surely safe for each fixed state.  It is
not strongly safe.

If $R_2=T\circ R_1$, then
$F_{z_1}^{R_1}\subseteq F_{T(z_1)}^{R_2}$ and consequently
\[
 \mathcal A_{R_2}(T(z_1))\subseteq\mathcal A_{R_1}(z_1).
\]
Thus deterministic coarsening can only shrink the common-action set.

\subsection{Least intervention on a common convex action set}

\begin{theorem}[Common-action projection]\label{thm:projection}
Suppose actions lie in $\mathbb R^p$ and every $\Usafe(x)$ in a fixed fiber is closed and convex.
If $\mathcal A_R(z)$ is nonempty, then for every nominal action $u_{\rm nom}(z)$ the problem
\begin{equation}
 \min_{u\in\mathcal A_R(z)}
 \frac12\|u-u_{\rm nom}(z)\|_2^2\label{eq:least-intervention}
\end{equation}
has a unique minimizer, denoted $\pi_{\rm LI}(z)$.  It is uniformly safe on the fiber and minimizes
intervention among all such actions.  Moreover,
\[
 \langle u_{\rm nom}(z)-\pi_{\rm LI}(z),u-\pi_{\rm LI}(z)\rangle\leq0
 \quad\forall u\in\mathcal A_R(z).
\]
\end{theorem}

\begin{proof}
The common set is a closed convex intersection.  A minimizing sequence for
\eqref{eq:least-intervention} has bounded objective and is therefore bounded.  In finite dimensions
it has a convergent subsequence; closedness makes the limit feasible and continuity makes it a
minimizer.  If two distinct minimizers $u,v$ existed, their feasible midpoint would satisfy
\[
 \left\|\frac{u+v}{2}-u_{\rm nom}\right\|^2
 =\frac12\|u-u_{\rm nom}\|^2+\frac12\|v-u_{\rm nom}\|^2
  -\frac14\|u-v\|^2,
\]
a strict improvement.  The variational inequality follows by taking the right derivative along
the feasible segment from the minimizer to any $u$ in the common set.
\end{proof}

The theorem is pointwise.  It does not ensure that $z\mapsto\pi_{\rm LI}(z)$ is measurable or
continuous, nor that infinitely many fiber constraints admit a finite certified reduction.  For
control-affine relative-degree-one barrier constraints the individual action sets are often
convex, but the barrier hypotheses needed for forward safety remain separate.

\subsection{Finite-action slack and randomization semantics}

Let the action set be $\{1,\ldots,k\}$ and let
$\ell(x,a)\in[0,+\infty)$ be a finite violation loss.  For one fiber $F=F_z^R$, define the
deterministic and pointwise expected-loss slacks, which may take extended-real values,
\[
 \beta_{\rm det}(z)=\min_a\sup_{x\in F}\ell(x,a),\qquad
 \beta_{\rm exp}(z)=\inf_{p\in\Delta_k}\sup_{x\in F}\sum_{a=1}^kp_a\ell(x,a).
\]
Here the state is worst-case after $p$ is fixed, but loss is averaged over the realized action.
This is weaker than strong/pathwise fiber safety.  The corresponding strong robust loss satisfies
\[
 \beta_{\rm strong}(z)
 =\inf_{p\in\Delta_k}\sum_{a=1}^kp_a\sup_{x\in F}\ell(x,a)
 =\beta_{\rm det}(z),
\]
because a linear objective on the simplex is minimized at a deterministic action.  Randomization
therefore gives no benefit under the strong semantics.

\begin{theorem}[Finite-action pointwise expected-loss slack]\label{thm:finite-action}
For every nonnegative loss,
\begin{equation}
 \frac1k\beta_{\rm det}(z)\leq\beta_{\rm exp}(z)\leq\beta_{\rm det}(z).
 \label{eq:finite-action-sandwich}
\end{equation}
If $\ell(x,a)=0$ exactly when $a\in\Usafe(x)$ and the fiber is action-conflicted, both slacks are
strictly positive.  For hard failure loss
$\ell(x,a)=\mathbf 1\{a\notin\Usafe(x)\}$,
\[
 \beta_{\rm det}(z)=\beta_{\rm strong}(z)=1,\qquad \beta_{\rm exp}(z)\geq 1/k.
\]
\end{theorem}

\begin{proof}
A Dirac mass on a deterministic minimizer gives the upper bound.  For any $p\in\Delta_k$, choose
$a_p$ with $p_{a_p}\geq1/k$.  Nonnegativity yields
\[
 \sup_x\sum_a p_a\ell(x,a)
 \geq p_{a_p}\sup_x\ell(x,a_p)
 \geq\frac1k\min_a\sup_x\ell(x,a).
\]
Taking the infimum over $p$ proves \eqref{eq:finite-action-sandwich}.  If the common set is empty,
every action is unsafe at some
compatible state.  Exact zero calibration makes each of the finitely many action-wise suprema
positive, and hard loss makes each equal to one.
\end{proof}

When $F=\{x_1,\ldots,x_n\}$ is finite and $L_{ia}=\ell(x_i,a)$, finite linear-program duality gives
the audit witness
\begin{equation}
 \beta_{\rm exp}(z)=\max_{q\in\Delta_n}\min_a\sum_{i=1}^n q_iL_{ia}.
 \label{eq:finite-fiber-dual}
\end{equation}
The primal minimizes $t\in\mathbb R$ subject to $Lp\leq t\mathbf1$ and $p\in\Delta_k$; its dual
maximizes $r\in\mathbb R$ subject to $L^Tq\geq r\mathbf1$ and $q\in\Delta_n$.

\subsection{Expected world-model loss does not control a worst fiber}

\begin{proposition}[Smooth localized collision]\label{prop:average-separation}
Fix $a\in(0,1)$, let $X$ be uniform on $[-1,1]$, set $h(x)=x$, and use identity dynamics.
For every $0<\delta<\min\{a,1-a\}$ there are smooth scalar encoder, decoder, and latent predictor
with
\[
 \epssep(R_\delta)\geq a,\qquad
 \mathbb E|X-D(R_\delta(X))|^2\leq2a^2\delta,
 \qquad P(R_\delta(x))=R_\delta(x).
\]
The same collided fiber is action-conflicted for actions $\{-1,+1\}$ and safe-action condition
$ux\geq0$.
\end{proposition}

\begin{proof}
Let
$\psi(t)=\exp(1-(1-t^2)^{-1})$ for $|t|<1$ and zero otherwise, and define
\[
 R_\delta(x)=x-a\psi((x-a)/\delta)+a\psi((x+a)/\delta).
\]
The two bump supports are disjoint and interior to $[-1,1]$, while
$R_\delta(a)=R_\delta(-a)=0$.  Thus a safe point of margin $a$ collides with an unsafe point.
With $D=P$ equal to the identity, reconstruction error is supported on intervals of total length
$4\delta$ and has magnitude at most $a$; the uniform density gives the stated bound.  Identity
dynamics makes latent prediction exact.  Finally, $a$ admits only action $+1$ and $-a$ only action
$-1$, so the common-action intersection is empty.
\end{proof}

Thus even smooth encoders with vanishing average reconstruction and prediction losses can retain
a fixed worst-case static and action obstruction.  Any positive implication needs tail coverage,
uniform regularity, or verification assumptions; a small empirical mean alone is insufficient.

\subsection{Limits of the theorem package}

First, $\epssep$ is $h$-relative: replacing $h$ by $ch$, $c>0$, preserves the safe set but scales
the defect by $c$.  Second, $R=h$ has zero static defect for one known scalar specification, yet
can merge equally safe states with disjoint safe-action sets; static separation alone gives
neither a nontrivial dimension lower bound nor safe-control sufficiency.  Third, Theorem
\ref{thm:exact-existence} optimizes over arbitrary latent subsets.  Its maximal set need not lie in
a continuous, smooth, neural, polynomial, or verifiable certificate class.

Positive continuous-action slack needs additional analytic control.  The disjoint closed sets
$A_1=\{(t,0):t\geq0\}$ and $A_2=\{(t,1/t):t>0\}$ in $\mathbb R^2$ have distance zero, so distance
loss has infimum robust slack zero despite an empty intersection.  Compactness or coercivity alone
is not a universal replacement for lower semicontinuity and uniform loss calibration.  Nonconvex
common sets can have nonunique projections, and nonclosed sets need not attain one.  Finally, every action result is
relative to the declared audited condition; common one-step feasibility does not establish latent
Markovianity, recursive feasibility, or long-horizon safety.  Stochastic encoders likewise require
a law-overlap notion not proved here, because Theorem~\ref{thm:data-processing} is an exact-fiber
result for deterministic maps.
